% ============================================================================
\section{Objet, sources et conventions}
% ============================================================================

\subsection{Objet et méthode}

\pr{PASAPAS} est la procédure générale de Cast3M pour les calculs incrémentaux : elle avance dans le temps et résout, à chaque instant, un problème mécanique non linéaire, un problème thermique, ou les deux.
Elle est écrite en Gibiane, le langage de commandes de Cast3M, et s'appuie sur une vingtaine de sous-procédures et sur des opérateurs écrits en Fortran (Esope).

Cette note décrit le code tel qu'il est écrit, et non tel que la documentation le présente.
Les procédures Gibiane ont été téléchargées sur \url{https://www-cast3m.cea.fr} (\texttt{download.php?nom=<procédure>\&type=procedure}), aux versions du tableau~\ref{tab:versions}.
Les opérateurs cités ont été vérifiés dans leur source Fortran, que le même site distribue (\texttt{download.php?nom=<routine>\&type=source}).
Les notices et les supports de formation n'ont servi qu'à orienter la lecture : ils peuvent être faux, et ils le sont sur au moins deux points (KTAN et CONF, § \ref{sec:ktan} et \ref{sec:conf}).
Les numéros de ligne cités renvoient à ces versions.
Rien de ce qui suit n'a été confirmé par un calcul : les constats du § \ref{sec:constats} méritent chacun un jeu de données d'essai.

\begin{table}[ht]
\centering\small
\caption{Procédures lues, avec la date de version portée en tête de fichier.}
\label{tab:versions}
\begin{tabular}{@{}lll@{}}
\toprule
Procédure & Version & Rôle\\
\midrule
\pr{PASAPAS}  & 26/06/09 & boucle en temps, couplage\\
\pr{UNPAS}    & 26/09/09 & un pas de mécanique (3\,843 lignes)\\
\pr{TRANSNON} & 26/03/20 & un pas de thermique non linéaire\\
\pr{TRANSLIN} & 24/09/11 & un pas de thermique linéaire\\
\pr{DUPONT2}  & 24/09/11 & un pas de thermique, schéma de Dupont\\
\pr{PAS\_DEFA} & 26/08/05 & valeurs par défaut\\
\pr{PAS\_MODL} & 26/05/28 & classification du modèle et du matériau\\
\pr{PAS\_INIT} & 26/10/15 & configurations initiales\\
\pr{PAS\_ETAT} & 26/04/14 & état aux instants intermédiaires\\
\pr{PAS\_RESU} & 25/03/07 & transfert et sauvegarde des résultats\\
\pr{CORMAN}   & 13/02/05 & méthode asymptotique numérique\\
\pr{UNILATER} & 26/04/21 & conditions unilatérales (contact)\\
\pr{PAS\_RAYO}, \pr{HRAYO}, \pr{HRCAV} & 25/05/21, 14/04/17, 12/06/07 & rayonnement\\
\pr{AUTOPILO}, \pr{PILOINDI}, \pr{PAS\_HELM}, \pr{DEFO\_IMP} & — & pilotage, Helmholtz, contrôle de déformation\\
\bottomrule
\end{tabular}
\end{table}

\subsection{Lire le pseudo-code}

Chaque algorithme est écrit en pseudo-code, dans un cadre intitulé « Algo~$n$ » (ou « Annexe~X » pour les options).
Les règles de lecture sont les suivantes.

\begin{itemize}
\item \textbf{Une ligne plus indentée s'exécute à l'intérieur de la ligne au-dessus} (corps d'une boucle ou d'un test). Il n'y a pas de « fin pour » ni de « fin si » : seule l'indentation porte l'imbrication.
\item \textbf{Deux lignes alignées s'exécutent l'une après l'autre.}
\item \textbf{« Algo~$n$ » ou « Annexe~X » dans une ligne est un appel} : le détail de cette étape est écrit ailleurs, sous ce nom.
\item $a \gets b$ est une affectation ; \Sortir{} quitte la boucle la plus intérieure ; \Retourner{} quitte l'algorithme.
\item Le commentaire grisé en bout de ligne donne les lignes du source, dans \pr{UNPAS} sauf mention contraire, ou le nom de la boucle Gibiane (\src{BEXTERN}, \src{ETIQ}…).
\end{itemize}

Les schémas ne contiennent aucune flèche entre boucles : une boucle contenue dans une autre est dessinée à l'intérieur de son cadre.

\subsection{Notations}

Les vecteurs et matrices discrets sont en gras : $\vc u$ (déplacements nodaux), $\vg\lambda$ (multiplicateurs de Lagrange), $\vc A$ (matrice des liaisons), $\vc K$ (rigidité), $\vc M$ (masse), $\vc C$ (amortissement, ou capacité thermique selon le contexte), $\vc F$ (forces).
Les champs aux points d'intégration sont aussi en gras : $\sig$ (contraintes), $\vg\eps$ (déformations), $\vc V$ (variables internes).
L'indice $n$ désigne le début du pas et $n+1$ sa fin ; l'exposant $(i)$ désigne l'itéré numéro $i$.
$\Delta$ désigne un incrément cumulé depuis le début du pas et $\delta$ la correction apportée par une itération.
$\norminf{\cdot}$ est le maximum des valeurs absolues sur tous les nœuds (ou points de Gauss) et toutes les composantes retenues.
« GD » abrège « grands déplacements » et « HPP » « hypothèse des petites perturbations ».

Les équations sont de deux natures.
Celles qui traduisent une ligne du source sont accompagnées de sa référence.
Les autres sont des rappels classiques, donnés pour expliquer ce que le code calcule ; le texte le signale.

\subsection{Trois règles de Gibiane utiles à la lecture}
\label{sec:gibiane}

\paragraph{Les expressions s'évaluent de gauche à droite, sans priorité d'opérateurs.}
C'est la règle la plus importante pour traduire le source en formules, et elle est établie par le Fortran de l'interpréteur :
\begin{itemize}
\item \texttt{ANASYN} découpe l'instruction en « phrases élémentaires » : il cherche la première parenthèse fermante et remonte à l'ouvrante correspondante. La parenthèse la plus intérieure est donc évaluée d'abord.
\item \texttt{PILOT} lit ensuite les opérateurs un à un, dans l'ordre où ils apparaissent dans la phrase.
\item Chaque opérateur prend ses opérandes dans une pile (\texttt{LIRABJ}) : le premier objet disponible du type voulu, en lisant de gauche à droite.
\item Le résultat est réinséré en tête de la partie non lue de la pile (\texttt{ECPI}), où l'opérateur suivant le trouve en premier.
\end{itemize}
Ainsi \texttt{1 + 2 * 3} vaut $(1+2)\times 3 = 9$, et l'expression \texttt{psc + 1.D0 ** 0.5D0} de \pr{CORMAN} vaut $(\mathrm{psc}+1)^{1/2}$.
Toutes les formules de cette note ont été dérivées en respectant cette règle.

\paragraph{La portée des noms est dynamique.}
À l'appel d'une procédure, \texttt{INIPRO} initialise chaque nom qu'elle emploie avec l'objet de même nom des niveaux appelants, s'il en existe un.
Une procédure peut donc lire une variable de son appelant sans la recevoir en argument (§ \ref{sec:translin}).

\paragraph{Lire un indice absent d'une table est une erreur.}
\texttt{REMPIL} lève l'erreur~535 quand on lit \texttt{T.'X'} et que l'indice \texttt{'X'} n'existe pas (§ \ref{sec:constats}).

% ============================================================================
\section{Vue d'ensemble}
\label{sec:vue}
% ============================================================================

\begin{glossaire}[Glossaire — objets partagés par toutes les procédures]
— & \src{PRECED} & Table passée à \pr{PASAPAS} par l'utilisateur : modèle, matériau, chargement, options, et résultats des appels précédents. Toutes les sous-procédures la reçoivent.\\
— & \src{WTAB} = \src{PRECED.'WTABLE'} & Table de travail, recréée par \pr{PAS\_DEFA} à chaque appel. Elle porte les options après défauts, les configurations, la raideur mémorisée, les compteurs.\\
— & \src{CONTINUATION} (\src{conti}) & État convergé au début du pas : temps $t_n$, $\vc u_n$, $\sig_n$, $\vc V_n$, $T_n$\ldots\\
— & \src{ESTIMATION} (\src{estim}) & État de fin de pas en construction ($t_{n+1}$, $\vc u_{n+1}$\ldots). \pr{PAS\_RESU} le recopie dans \src{CONTINUATION} en fin de pas.\\
$\{t_k\}$ & \src{TEMPS_CALCULES} & Instants à atteindre, dans l'ordre.\\
— & \src{TEMPS_SAUVES} & Instants dont les résultats sont sauvegardés dans les tables de sortie.\\
— & \src{MODELE}, \src{CARACTERISTIQUES} & Modèle (maillage, formulation, loi) et matériau ; ils peuvent évoluer dans le temps (\src{CHAR_MODE}).\\
— & \src{CHARGEMENT} & Chargements tabulés en temps (\op{TIRE} les évalue à un instant).\\
— & \src{FOR0}, \src{GE0_DEB}, \src{FOR} & Configurations : de référence, de début de pas, de fin de pas (§ \ref{sec:configs}).\\
\end{glossaire}

\subsection{Le problème résolu}

\paragraph{Mécanique.}
Après discrétisation par éléments finis, on cherche à chaque instant $t_{n+1}$ les déplacements $\vc u_{n+1}$ et les multiplicateurs $\vg\lambda_{n+1}$ tels que
\begin{equation}
\label{eq:equilibre}
\vc F^{int}(\vc u_{n+1}) + \vc A^{\mathsf T}\vg\lambda_{n+1} = \vc F^{ext}(t_{n+1}) + \vc F^{suiv}(\vc u_{n+1}),
\qquad
\vc A\,\vc u_{n+1} = \vc u^{imp}(t_{n+1}),
\end{equation}
où $\vc F^{int} = \int_\Omega \vc B^{\mathsf T}\sig\,\mathrm d\Omega$ (opérateur \op{BSIGMA}), $\sig$ étant obtenue en intégrant la loi de comportement sur le pas, $\vc F^{suiv}$ regroupe les chargements suiveurs (pression en GD, procédure utilisateur \src{CHARMECA}), et $\vc A\vc u = \vc u^{imp}$ représente les déplacements imposés et les relations.
Les blocages sont traités par multiplicateurs de Lagrange : $-\vc A^{\mathsf T}\vg\lambda$ est la réaction.
En dynamique, on ajoute $\vc M\ddot{\vc u} + \vc C\dot{\vc u}$ au premier membre ; en consolidation, un terme de couplage avec la pression de pore.

\paragraph{Thermique.}
On cherche les températures nodales $\vc T$ telles que
\begin{equation}
\label{eq:thermique}
\vc C(\vc T)\,\dot{\vc T} + \vc K(\vc T)\,\vc T = \vc F(\vc T, t),
\end{equation}
où $\vc C$ est la capacité, $\vc K$ regroupe conduction, convection, advection et rayonnement linéarisé, et $\vc F$ les flux, sources, températures extérieures de convection et le second membre du rayonnement ; les températures imposées passent aussi par des multiplicateurs.

\paragraph{Couplage.}
La thermique agit sur la mécanique par le matériau, évalué à la température courante, et par la déformation thermique $\vg\eps^{th} = \op{EPTH}(T) - \op{EPTH}(T_{ref})$.
La mécanique n'agit sur la thermique que par la géométrie en GD et par les procédures utilisateur (§ \ref{sec:couplage}).

\subsection{Structure générale}

La figure~\ref{fig:imbrication} montre l'imbrication des boucles : chaque cadre est une boucle, contenue dans celle qui l'entoure.
Dans une passe, la thermique (cadre de gauche) s'exécute entièrement avant la mécanique (cadre de droite).

\begin{figure}[ht]
\centering
\begin{tikzpicture}[x=1cm,y=1cm,
  bx/.style={draw=algframe,rounded corners=2pt,line width=0.5pt},
  lb/.style={anchor=north west,font=\sffamily\scriptsize,inner sep=2pt,align=left}]
  \draw[bx,fill=algtitle!35] (0,0) rectangle (16,10.2);
  \node[lb,text width=15.6cm] at (0.05,10.15) {\textbf{PASAPAS} — \algoref{1}};
  \draw[bx,fill=white] (0.3,0.3) rectangle (15.7,9.55);
  \node[lb,text width=15cm] at (0.35,9.5) {pour chaque instant $T_I$ de \src{TEMPS_CALCULES} — boucle \src{BEXTERN}};
  \draw[bx,fill=algtitle!20] (0.6,0.6) rectangle (15.4,8.85);
  \node[lb,text width=14.5cm] at (0.65,8.8) {pour chaque passe $k$, une seule par défaut — boucle \src{BO_BOTH}};
  % thermique
  \draw[bx,draw=ther,fill=ther!6] (0.9,0.9) rectangle (6.7,8.1);
  \node[lb,text width=5.6cm] at (0.95,8.05) {\textbf{\textcircled{\tiny 1} thermique de $t$ à $T_I$}\\ \algoref{2}, \algoref{2L} ou \algoref{2D}};
  \draw[bx,draw=ther,fill=white] (1.2,1.2) rectangle (6.4,6.95);
  \node[lb,text width=5cm] at (1.25,6.9) {pour chaque sous-pas thermique\\ boucle \src{BOUC_SO} (si \src{PAS_MAX})};
  \draw[bx,draw=ther,fill=ther!6] (1.5,1.5) rectangle (6.1,5.8);
  \node[lb,text width=4.5cm] at (1.55,5.75) {répéter les itérations — \src{BOU_IT1}\\ (\pr{TRANSNON} seulement)};
  \node[font=\sffamily\scriptsize,align=center,text width=4.2cm] at (3.8,3.4) {$\vc S\,\Delta\vc T^{(i)} = \vc R^{(i-1)}$\\[2pt] critère sur la variation\\ de $\Delta\vc T$};
  % mécanique
  \draw[bx,draw=meca,fill=meca!6] (7,0.9) rectangle (15.1,8.1);
  \node[lb,text width=7.9cm] at (7.05,8.05) {\textbf{\textcircled{\tiny 2} mécanique de $t$ à $T_I$} — \pr{UNPAS}, \algoref{3}};
  \draw[bx,draw=meca,fill=white] (7.3,1.2) rectangle (14.8,7.35);
  \node[lb,text width=7.3cm] at (7.35,7.3) {pour chaque sous-pas $s$, 100 au plus — \src{BONOCONV} (convergence forcée)};
  \draw[bx,draw=meca,fill=meca!6] (7.6,1.5) rectangle (14.5,6.6);
  \node[lb,text width=6.8cm] at (7.65,6.55) {répéter les itérations d'équilibre $i$ — \src{ETIQ}, \algoref{4} ; tests : \algoref{6}};
  \draw[bx,draw=meca,fill=white] (7.9,1.8) rectangle (14.2,5.85);
  \node[lb,text width=6.2cm] at (7.95,5.8) {pour chaque sous-incrément $k$ — \src{sousinc}, \algoref{5}};
  \draw[bx,draw=meca,fill=meca!6] (8.2,2.1) rectangle (13.9,5.1);
  \node[lb,text width=5.6cm] at (8.25,5.05) {répéter les sous-itérations non locales — \src{BIHL}, \annref{A} (si \src{NON_LOCAL})};
  \node[font=\sffamily\scriptsize,align=center,text width=5cm] at (11.05,3.2) {loi de comportement \op{COMP}\\ puis régularisation de $\vc V$};
\end{tikzpicture}
\caption{Imbrication des boucles. Un cadre contenu dans un autre est une boucle intérieure. Dans une passe, le bloc \textcircled{\tiny 1} s'exécute avant le bloc \textcircled{\tiny 2}.}
\label{fig:imbrication}
\end{figure}

\begin{algobox}{0}{Imbrication des boucles}{PASAPAS · TRANSNON · UNPAS}
\begin{algorithmic}[1]
\Etape{\pr{PASAPAS}} \Comment{\algoref{1}}
\ForAll{instant $T_I$} \Comment{boucle \src{BEXTERN}}
  \ForAll{passe $k$, une seule par défaut} \Comment{boucle \src{BO_BOTH}}
    \Etape{thermique de $t$ à $T_I$} \Comment{\algoref{2}, \algoref{2L} ou \algoref{2D}}
      \ForAll{sous-pas thermique} \Comment{\src{BOUC_SO}, si \src{PAS_MAX}}
        \State répéter les itérations \Comment{\src{BOU_IT1}}
      \EndFor
    \EndEtape
    \Etape{puis mécanique de $t$ à $T_I$} \Comment{\algoref{3}, \pr{UNPAS}}
      \ForAll{sous-pas $s$, 100 au plus} \Comment{\src{BONOCONV}, convergence forcée}
        \Etape{répéter les itérations $i$} \Comment{\src{ETIQ}, \algoref{4}}
          \ForAll{sous-incrément $k$} \Comment{\src{sousinc}, \algoref{5}}
            \State répéter les sous-itérations non locales \Comment{\src{BIHL}, \annref{A}}
          \EndFor
        \EndEtape
      \EndFor
    \EndEtape
  \EndFor
\EndFor
\EndEtape
\end{algorithmic}
\end{algobox}

\subsection{L'essentiel en six points}

\begin{enumerate}
\item \textbf{Boucle en temps.} On parcourt une liste d'instants imposés (\src{TEMPS_CALCULES}). Si l'on ne donne que \src{TEMPS_SAUVES}, le pas s'adapte au nombre d'itérations : divisé par $\sqrt2$ au-delà de 20, multiplié par $\sqrt2$ en deçà de 7 (§ \ref{sec:temps}).
\item \textbf{Couplage étagé.} La thermique avance de $t_n$ à $t_{n+1}$, puis la mécanique utilise $T_{n+1}$. Il n'y a qu'une passe par défaut (§ \ref{sec:couplage}).
\item \textbf{Newton modifié, accéléré.} L'opérateur des itérations est par défaut la rigidité élastique (plus $\vc K_\sigma$ en GD), gardée d'un pas à l'autre quand c'est possible. L'accélération \op{ACT3} intervient une itération sur deux. La tangente n'est utilisée qu'avec \src{K_TANGENT}, et ce n'est pas la tangente algorithmique (§ \ref{sec:operateur}).
\item \textbf{Deux critères à $10^{-4}$.} Le résidu d'équilibre, en norme max et rapporté à une force de référence fixée en début de pas, et la variation de l'incrément de déformation entre deux itérations ; un troisième critère porte sur les moments (§ \ref{sec:criteres}).
\item \textbf{Échec.} La non-convergence est détectée quand le critère n'a pas baissé de 1\,\% en 10 itérations, ou après \src{MAXITERATION} itérations. Par défaut, on repart alors de l'itéré non convergé comme nouvel état initial (convergence forcée), jusqu'à 100 sous-pas (§ \ref{sec:pas}).
\item \textbf{Configurations.} En HPP, tout se fait sur la configuration initiale. En GD : début de pas pour la rigidité initiale et $\vc F^{int}_n$, mi-incrément pour $\Deps$ et la loi de comportement, fin de pas recalculée à chaque itération pour les forces internes (§ \ref{sec:configs}).
\end{enumerate}
